\documentclass[10pt,a4paper]{amsart}
\usepackage[margin=19mm]{geometry}
\usepackage{amsmath,amssymb,mathtools}
\usepackage[scr=rsfs,cal=euler]{mathalfa}
\usepackage{xfrac}
\usepackage[unicode,hidelinks,bookmarks=false]{hyperref}
\newcommand{\ie}{i.e.\ }
\newcommand{\cf}{cf.\ }
\newcommand{\resp}{respectively\ }
\newcommand{\Subsection}{\subsection}
\providecommand{\nobreakdash}{\nobreak\hbox{-}}
\setlength{\emergencystretch}{2em}
\multlinegap0pt
\usepackage{amssymb}
\usepackage{mathtools}
\usepackage{xcolor}
\usepackage{enumerate}
\newtheorem{prop}{Proposition}
\renewcommand{\P}{\mathbb{P}}
\newcommand{\Z}{\mathbb{Z}}
\title{Bicomplex Characteristic Classes}
\author{Jos\'e Antonio Arciniega-Nev\'arez, Yesenia Bravo-Ortega, In\'acio Rabelo, Agust\'in Romano-Vel\'azquez}
\date{Source: arXiv:2609.17692v1 (CC BY 4.0)}
\begin{document}
\maketitle

\noindent\emph{Source and licence.} Bicomplex Characteristic Classes. Version arXiv:2609.17692v1 (CC BY 4.0), distributed by arXiv under the Creative Commons Attribution 4.0 International licence, http://creativecommons.org/licenses/by/4.0/, as recorded in the arXiv licence metadata for this exact version and shown on the abstract page https://arxiv.org/abs/2609.17692v1. Full licence notice: \texttt{licenses/arxiv-2609.17692-CC-BY-4.0.txt}.\par\smallskip

\noindent\textit{Editorial scope.} Counted unit: the article's prop prop:splitting (source0.tex lines 880--887). The article's complete original proof of it is delivered with it (source0.tex lines 889--912, one \texttt{proof} environment of 24 lines). No mathematics is written, repaired or re-derived: the statement and the proof are the article's own lines, in the article's own notation, and no retained line is substituted or re-flowed.  No further source-local context is needed: every cross-reference of every retained line resolves inside the two delivered spans.  Nothing was omitted from any retained span, so the package carries no omission record.  No retained line contains a citation, so the wrapper carries no bibliography and no bibliography entry.  No retained line contains a figure environment, an image-inclusion command or drawing code, so no image is delivered and the package declares no asset dependency.  Editorial formatting only, in addition: an AMS \texttt{amsart} preamble that restates the article's own declarations and macros used by the retained lines, namely the environments \texttt{prop} on separate counters, together with the standard packages those lines need.  The retained environments print in this document as Proposition 1 in order of appearance, whereas the article prints the same items with the numbering of its own full document.  The complete native source of the article as distributed in the arXiv e-print for this exact version is bundled unchanged as \texttt{sources/210620/source0.tex}.
\begin{prop}[Splitting principle]\label{prop:splitting}
    Let $\pi:E \to B$ be a bicomplex fiber bundle. Then there exists a continuous map $f: P \to B$ such that
    \begin{enumerate}
        \item $f^*E$ is a direct sum of bicomplex line bundles on $P$;
        \item $f^*: H^*(B;\Z) \to H^*(P;\Z)$ is injective.
    \end{enumerate}
    The map $f$ is called a \emph{splitting map}.
\end{prop}

\begin{proof}
    The proof is identical to the complex case. We proceed by induction on the rank $n$. For $n=1$, take $P=B$ and $f=\mathrm{id}$. Assume the result holds for ranks $< n$. Let $\P\pi: \P E \to B$ be the bicomplex projective bundle. Over $\P E$, we have the tautological line subbundle $L(\pi) \subset \pi^*E$. Consider the idempotent decomposition
\begin{equation*}
\pi^*E=\pi^*E^1e\oplus\pi^*E^2e^\dagger.
\end{equation*}
The tautological line bundle decomposes accordingly as
\begin{equation*}
L(\pi)=L^1(\pi)e\oplus L^2(\pi)e^\dagger.
\end{equation*}
Taking the orthogonal complements of $L^1(\pi)$ in $\pi^*E^1$ and of $L^2(\pi)$ in $\pi^*E^2$, we obtain complex subbundles $Q^1$ and $Q^2$. Hence,
\begin{equation*}
Q=Q^1e\oplus Q^2e^\dagger
\end{equation*}
is a bicomplex subbundle, and we obtain the splitting sequence
\begin{equation*}
0\longrightarrow L(\pi)\longrightarrow \pi^*E\longrightarrow Q\longrightarrow 0,
\end{equation*}
Moreover,
\begin{equation*}
\pi^*E=L(\pi)\oplus Q.
\end{equation*}

The bundle $Q$ has rank $n-1$. By the induction hypothesis, there exists a map $g: P' \to \P E$ such that $g^*Q$ splits into line bundles and $g^*$ is injective on cohomology. Let $f = \P\pi \circ g: P' \to B$. Then $f^*E = g^*\pi^*E \cong g^*L(\pi) \oplus g^*Q$, where $g^*L(\pi)$ is a bicomplex line bundle and $g^*Q$ splits by induction. The injectivity of $f^*$ follows from the injectivity of $\P\pi^*$ (Leray--Hirsch) and of $g^*$.
\end{proof}

\end{document}
